\documentclass[10pt,a4paper]{article}
\usepackage[utf8]{inputenc}
\usepackage[T1]{fontenc}
\usepackage{lmodern}
\usepackage[french]{babel}
\usepackage[margin=1.6cm]{geometry}
\usepackage{microtype}
\usepackage{amsmath,amssymb}
\usepackage{booktabs}
\usepackage{array}
\usepackage[table]{xcolor}
\usepackage{enumitem}
\usepackage[most]{tcolorbox}
\definecolor{navy}{HTML}{1F3864}
\definecolor{bluemid}{HTML}{2E5496}
\definecolor{lightblue}{HTML}{D6E0F0}
\definecolor{accent}{HTML}{C0491F}
\newtcolorbox{cle}[1][]{colback=lightblue!40,colframe=navy,boxrule=0.6pt,
  arc=2pt,left=6pt,right=6pt,top=4pt,bottom=4pt,before skip=5pt,after skip=5pt,#1}
\setlength{\parindent}{0pt}
\setlength{\parskip}{3pt}
\pagestyle{empty}

\begin{document}
\begin{center}
{\color{navy}\Large\bfseries Quantifier le capital d'un défaut de conformité DORA}\\[2pt]
{\color{bluemid}\large Une cascade dirigée entre les cinq piliers, et la frontière de ce qui est mesurable}\\[3pt]
{\small Kélian Kaddouri \textbullet{} ENSAE / Nexialog Consulting \textbullet{} résumé exécutif}
\end{center}
\vspace{-2pt}

\begin{cle}
\textbf{La thèse en une phrase.} La dépendance entre les cinq piliers de DORA n'est pas une
corrélation mais une \textbf{propagation orientée}~; or cette orientation \emph{n'est pas
mesurable} sur les données publiques. Plutôt que de la poser en silence, ce mémoire en établit
la \textbf{frontière d'identifiabilité}, publie le capital en \textbf{bornes} sur tout ce que la
donnée laisse ouvert, et transforme la donnée manquante en \textbf{spécification de reporting
chiffrée et hiérarchisée}.
\end{cle}

\textbf{Le modèle.} Un incident amorcé sur un pilier se propage aux autres selon une matrice
d'influence \emph{asymétrique} $W$, normalisée à la Leontief : la stabilité du système est
\textbf{garantie par construction}, non supposée. La sévérité est une Pareto généralisée (théorie
des valeurs extrêmes), la fréquence une négative binomiale, et l'état de conformité de l'entité
pilote le capital par une couche multi-états markovienne, d'où une \textbf{trajectoire} de
capital plutôt qu'un chiffre. Là où la Formule Standard, qui ne dépend que du volume, est
structurellement \textbf{aveugle} à la conformité.

\textbf{La contribution centrale : une frontière, pas une mesure.} Sur deux sources qui échouent
pour des raisons \emph{opposées} (l'une manque de taxonomie, l'autre de résolution temporelle),
et jusque sur la meilleure expérience naturelle disponible (le choc de chaîne
d'approvisionnement MOVEit, où l'effet apparent est intégralement absorbé par un placebo), on
démontre que la \textbf{direction} de la contagion n'est pas identifiable, quand sa
\textbf{co-occurrence} l'est. La théorie des processus réversibles en donne la raison de
principe : la direction est un \emph{courant} qui change de signe sous renversement du temps,
et aucune statistique invariante par ce renversement ne peut l'identifier.

\begin{cle}
\textbf{Ce que la donnée fixe, ce qu'elle ne fixe pas, et ce qu'on en fait.}
\begin{itemize}[leftmargin=12pt,itemsep=1pt,topsep=2pt]
\item \textbf{Fixé} et validé : la queue de sévérité (ajustement non rejeté, $p=0{,}87$ par
      bootstrap paramétrique), la surdispersion de la fréquence ($p\approx10^{-30}$), la
      co-occurrence entre piliers, l'accumulation par prestataire tiers.
\item \textbf{Non fixé} : la direction. Elle est donc \textbf{bornée}, jamais posée, et le
      capital est lu sur l'ensemble entier des matrices admissibles.
\item \textbf{Corroboré} sans élicitation : le codage de sept post-mortems officiels (CSRB,
      FCA/PRA, GAO, SEC, OCC) donne une orientation \emph{cohérente} d'un rapport à l'autre, et
      un traitement bayésien conjugué en tire une loi a posteriori qui \emph{reproduit d'elle-même}
      la structure d'identification, robuste à un prior sceptique ($0{,}8\,\%$ d'écart).
\end{itemize}
\end{cle}

\textbf{Les résultats qui résistent.} Le \textbf{niveau} du capital est illustratif : il bouge
d'un facteur $40$ selon la source de sévérité, d'un facteur $5{,}5$ selon la famille de loi de
queue, et d'un facteur $2{,}6$ par la seule incertitude d'estimation. Ce qui \textbf{résiste} est
ailleurs, et c'est ce que le mémoire défend : le \textbf{classement des piliers} (gouvernance et
gestion des tiers dominent robustement les trois autres, dans tous les cadres de dépendance
testés), les \textbf{rapports} entre états de conformité, et l'\textbf{ordre de remédiation}.
Trois étages d'incertitude sont nommés et chiffrés séparément : paramètre, identification,
et \textbf{modèle} (le choix de la famille elle-même) --- ce dernier étage étant celui que la
plupart des travaux passent sous silence.

\textbf{La portée réglementaire.} La limite devient un livrable. Le registre d'incidents qui
identifierait la contagion est spécifié (horodatage de survenance au mois, taxonomie par domaine
de contrôle, causes communes consolidées), et surtout \textbf{hiérarchisé} : documenter une seule
dépendance inter-piliers lève jusqu'à $29\,\%$ de l'ambiguïté de capital, une autre n'en lève
rien. La recommandation cesse d'être générique. En regard, chaque indicateur observable
(délai de notification, durée de confinement, concentration tierce) est rattaché à un canal du
modèle et devient un \textbf{levier de capital priorisable}, ce qui fait de la conformité un
arbitrage économique et non un pur coût.

\begin{cle}[colback=accent!7,colframe=accent!70!black]
\textbf{Ce que ce mémoire ne prétend pas.} Le niveau absolu du capital n'est pas une mesure et
n'est jamais présenté comme telle : l'agrégat obtenu est d'échelle sectorielle, non celui d'une
entité isolée. La direction de la contagion n'est pas calibrée, seulement bornée et corroborée,
avec son confondant déclaré (une enquête officielle cherche par convention une cause racine de
gouvernance, ce qui peut orienter la lecture). Connaître le sens d'un lien ne donnerait d'ailleurs
pas sa force : même les dix directions établies, une bande subsisterait. \textbf{L'identification
partielle n'est donc pas un pis-aller provisoire, elle est structurelle} --- et c'est précisément
ce que le mémoire assume, chiffre et documente, plutôt que de le masquer derrière un nombre unique.
\end{cle}

\vspace{-1pt}
{\footnotesize\textcolor{navy}{\textbf{Mots-clés.}} Risque cyber \textbullet{} DORA \textbullet{}
SCR \textbullet{} Solvabilité~II \textbullet{} valeurs extrêmes \textbullet{} cascade dirigée
\textbullet{} identification partielle \textbullet{} bornes de capital \textbullet{} inférence
bayésienne \textbullet{} valeur de l'information.}

\end{document}
